% =====================================================================
\chapter{Выходы и кривые}
% =====================================================================
\chapterintro
  {как настроить направление сервомашинок, пределы хода и нейтраль; что такое кривые и
   как построить свою кривую газа}
  {модель с сервами (для практики), либо полётный контроллер с вкладкой Receiver}

\section{Страница Outputs}

\mpath{MDL\sep Outputs} --- последний этап обработки канала перед радиомодулем.
Здесь настраивается \emph{механика}: как канал ложится на конкретную серву.

\begin{center}
\lcd{\inv{OUTPUTS}\\
CH1\quad 0.0\quad -100\quad 100\quad ---\\
CH2\quad 0.0\quad -100\quad 100\quad INV\\
CH3\quad 0.0\quad -100\quad 100\quad ---\\
CH4\quad 2.5\quad \hphantom{0}-85\quad \hphantom{0}85\quad ---}
\end{center}

\begin{center}\small
\renewcommand{\arraystretch}{1.2}
\begin{tabularx}{\linewidth}{@{}p{30mm}X@{}}
\toprule
\menu{Name} & имя канала (удобно: \code{AilL}, \code{Gear}) \\
\menu{Subtrim} & небольшой сдвиг нейтрали сервы, $\pm100\,\%$ \\
\menu{Min}, \menu{Max} & ограничители хода. Канал \emph{никогда} не выйдет за эти значения,
  что бы ни делали миксеры \\
\menu{Direction} & \menu{INV} --- инверсия (реверс сервы) \\
\menu{Curve} & кривая на выходе (редко нужна) \\
\menu{PPM Center} & центр импульса в мкс (по умолчанию 1500) --- смещает всё, включая пределы \\
\menu{Subtrim mode} & \menu{$\Delta$} --- сабтрим сдвигает весь ход; \menu{=} --- симметричный:
  пределы остаются на месте \\
\bottomrule
\end{tabularx}
\end{center}

\section{Порядок настройки серв}

\begin{enumerate}
  \item \textbf{Механически} выставьте качалку сервы максимально близко к нейтрали
        (переставьте её на шлицах). Сабтрим --- только для доводки.
  \item \textbf{Направление}: отклоните стик и проверьте, куда пошла поверхность. Неверно ---
        \menu{Direction} = \menu{INV}. Реверсируйте именно здесь, а не весом во входах,
        иначе перепутаются знаки в миксерах.
  \item \textbf{Нейтраль}: \menu{Subtrim}, пока поверхность не встанет ровно.
  \item \textbf{Пределы}: уменьшайте \menu{Min}/\menu{Max}, пока поверхность при полном
        отклонении стика не перестанет упираться в механику (серва не должна «жужжать»
        под нагрузкой в крайнем положении).
\end{enumerate}

\begin{tip}
Долгое \btn{ENTER} на строке выхода даёт команды \menu{Trims to subtrims} (перенести
значения триммеров в сабтримы и обнулить триммеры) и \menu{Copy stick to subtrim}. Удобно
после облёта: слетали, оттриммировали, перенесли --- и триммеры снова в нуле.
\end{tip}

\begin{caution}
Для полётного контроллера (Betaflight, INAV, ArduPilot) каналы \sw{CH1}--\sw{CH4} должны давать
ровно 1000/1500/2000\,мкс или близко к этому. Не ставьте на них сабтримы и пределы;
проверяйте значения на вкладке \menu{Receiver} конфигуратора контроллера. Если крайние
значения 988/2012\,мкс --- это нормально (±100\,\% в EdgeTX), контроллер всё равно поймёт.
\end{caution}

\section{Кривые}

\textbf{Кривая} --- произвольная функция $y = f(x)$, заданная точками. Используются
во входах, миксерах и выходах. \mpath{MDL\sep Curves}: до 32 кривых, в каждой
от 2 до 17 точек.

\begin{itemize}
  \item \menu{Type}: \menu{Standard} --- точки равномерно по $x$; \menu{Custom} --- можно
        двигать и $x$, и $y$.
  \item \menu{Count} --- число точек.
  \item \menu{Smooth} --- сглаживание (кривая проходит плавно через точки).
\end{itemize}

\subsection{Пример: кривая газа}

Для самолёта с ДВС или для плавного висения удобна кривая, которая даёт мягкий газ на
малых оборотах и быстрый набор вверху.

\begin{figure}[H]
\centering
\begin{tikzpicture}
\begin{axis}[width=9cm,height=6.2cm,xmin=-100,xmax=100,ymin=-100,ymax=100,
  xlabel={стик газа, \%},ylabel={выход, \%},grid=major,grid style={gray!20},
  tick label style={font=\scriptsize},label style={font=\small},
  legend pos=north west,legend style={font=\scriptsize}]
\addplot[thick,gray,domain=-100:100] {x};
\addplot[thick,etx,mark=*,mark options={fill=white}] coordinates {(-100,-100) (-50,-70) (0,-25) (50,35) (100,100)};
\addplot[thick,accent,smooth,dashed] coordinates {(-100,-100) (-50,-70) (0,-25) (50,35) (100,100)};
\legend{линейно,5 точек,5 точек + Smooth}
\end{axis}
\end{tikzpicture}
\caption{Кривая газа из пяти точек: $-100, -70, -25, 35, 100$}
\end{figure}

\subsection{Кривая в миксере или во входе?}

Кривую во \emph{входе} видит «пилот»: она меняет ощущение стика во всех каналах, куда этот
вход попадает. Кривую в \emph{миксере} --- только конкретный канал. Кривую на \emph{выходе}
применяют для компенсации нелинейной механики (например, кривошипной тяги).

\subsection{Функции вместо кривых}

Вместо своей кривой можно выбрать встроенную функцию: \menu{x>0}, \menu{x<0}, \menu{|x|},
\menu{f>0}, \menu{f<0}, \menu{|f|}. Например, функция \menu{x>0} на крутилке даёт
выход только во второй половине её хода.

\begin{tasks}
\begin{enumerate}
  \item Подключите серву к приёмнику на \sw{CH1}, ограничьте её ход до $\pm80\,\%$ и
        убедитесь, что крен стиком больше не выводит серву за этот предел.
  \item Создайте кривую из 5 точек и назначьте её на газ во входе \sw{I3}. Проследите
        значения на экране каналов.
\end{enumerate}
\end{tasks}
